\documentclass[10pt,a4paper]{article}
\usepackage[a4paper,margin=1.55cm,top=1.4cm,bottom=1.55cm]{geometry}
\usepackage{iftex}
\ifPDFTeX
  \usepackage[T1]{fontenc}
  \usepackage[utf8]{inputenc}
  \usepackage{lmodern}
\else
  \usepackage{fontspec}
  \setmainfont{DejaVu Serif}
  \setsansfont{DejaVu Sans}
\fi
\usepackage{microtype}
\usepackage{xcolor}
\usepackage{enumitem}
\usepackage{titlesec}
\usepackage{tabularx}
\usepackage{array}
\usepackage{hyperref}
\usepackage{fancyhdr}
\usepackage{lastpage}
\usepackage{ragged2e}
\usepackage{amsmath}
\usepackage{textcomp}
\usepackage{needspace}


% ===============================================================
% ===============================================================
% ===============================================================
% Muted brown-and-green academic palette
% Forest green: name, section headings, and structural emphasis
\definecolor{navy}{HTML}{365B4A}
% Walnut brown: professional title and hyperlinks
\definecolor{teal}{HTML}{795A43}
% Deep taupe: secondary structural text
\definecolor{slate}{HTML}{61564C}
% Pale sage: understated rules and dividers
\definecolor{rulegray}{HTML}{C8CEC4}
\hypersetup{colorlinks=true,urlcolor=teal,linkcolor=navy,pdfauthor={Vigneshwaran Krishnamurthy},pdftitle={Curriculum Vitae -- Vigneshwaran Krishnamurthy}}
\setlength{\parindent}{0pt}
\setlength{\parskip}{2.5pt}
\emergencystretch=1em
\clubpenalty=10000
\widowpenalty=10000
\pretocmd{\section}{\Needspace{6\baselineskip}}{}{}
\setlist[itemize]{leftmargin=1.25em,itemsep=1.5pt,topsep=2pt,parsep=0pt}
\titleformat{\section}{\large\sffamily\bfseries\color{navy}}{}{0pt}{}[\color{rulegray}\titlerule]
\titlespacing*{\section}{0pt}{9pt}{4pt}
%\pagestyle{fancy}
%\fancyhf{}
%\fancyhead[L]{\sffamily\small\color{slate}Vigneshwaran Krishnamurthy}
%\fancyhead[R]{\sffamily\small\color{slate}Curriculum Vitae}
%\fancyfoot[C]{\sffamily\scriptsize\color{slate}\thepage\ / \pageref*{LastPage}}
%\renewcommand{\headrulewidth}{0.3pt}

% ===============================================================
% ===============================================================
% ===============================================================

\newcommand{\cventry}[5]{%
  \begin{tabularx}{\textwidth}{@{}>{\raggedleft\arraybackslash}p{2.55cm}@{\hspace{0.55cm}}X@{}}
  \sffamily\small\color{slate}#1 & \textbf{#2}\\
  & \emph{#3}\if\relax\detokenize{#4}\relax\else, #4\fi\\
  & #5
  \end{tabularx}\par\vspace{3pt}}

% Lighter color for presentation titles, hosts, and dates
% Soft olive-brown: descriptions, dates, hosts, and bullet points
\definecolor{entrylight}{HTML}{746B58}

% Presentation entry:
% #1 Conference and presentation type
% #2 Location
% #3 Presentation title
% #4 Date
\newcommand{\presentationentry}[4]{%
\begin{tabularx}{\textwidth}{
    @{}
    >{\raggedright\arraybackslash}X
    >{\raggedleft\arraybackslash}p{4.0cm}
    @{}
}
#1 & #2 \\
\textcolor{entrylight}{#3} &
\textcolor{entrylight}{\emph{#4}}
\end{tabularx}
\vspace{2pt}
}

% ===============================================================

% Invited-talk entry:
% #1 Institution or laboratory
% #2 Location
% #3 Host
% #4 Date
\newcommand{\invitedtalkentry}[4]{%
\begin{tabularx}{\textwidth}{
    @{}
    >{\raggedright\arraybackslash}X
    >{\raggedleft\arraybackslash}p{4.5cm}
    @{}
}
#1 & #2 \\
\textcolor{entrylight}{Host: #3} &
\textcolor{entrylight}{\emph{#4}}
\end{tabularx}
\vspace{2pt}
}

% ===============================================================

% Grant/award heading: title on left, date on right
\newcommand{\grantheading}[2]{%
\begin{tabularx}{\textwidth}{
    @{}
    >{\raggedright\arraybackslash}X
    >{\raggedleft\arraybackslash}p{4.5cm}
    @{}
}
\textbf{#1} & #2
\end{tabularx}
}

% Lighter-colored bullet points
\newenvironment{lightitems}
{%
  \begin{itemize}[
      label=\textcolor{entrylight}{\textbullet},
      leftmargin=1.3em,
      itemsep=1pt,
      topsep=2pt,
      parsep=0pt
  ]
  \color{entrylight}
}
{%
  \end{itemize}
}

% ===============================================================

% Project entry:
% #1 Project title
% #2 Date
% #3 Role
% #4 Description
\newcommand{\projectentry}[4]{%
\begin{tabularx}{\textwidth}{
    @{}
    >{\raggedright\arraybackslash}X
    >{\raggedleft\arraybackslash}p{2.5cm}
    @{}
}
\textbf{#1} & #2
\end{tabularx}

#3

{\color{entrylight}#4\par}

\vspace{4pt}
}

% ===============================================================
% ===============================================================
% ===============================================================
% Publication records checked on 24 November 2026.
\begin{document}
\thispagestyle{empty}
{\fontsize{25}{28}\selectfont\sffamily\bfseries\color{navy} Vigneshwaran Krishnamurthy}\par
\vspace{2pt}
{\large\sffamily\color{teal}Specially Appointed Assistant Professor}\par
{\small Astronomical Institute, Graduate School of Science, Tohoku University $\cdot$ Sendai, Japan}\par
\vspace{4pt}
{\small\sffamily
\href{mailto:vignesh.krishnamurthy@astr.tohoku.ac.jp}{vignesh.krishnamurthy@astr.tohoku.ac.jp}
\quad$\vert$\quad
\href{https://vigneshastro.github.io/}{Website}
\quad$\vert$\quad
\href{https://orcid.org/0000-0003-2310-9415}{ORCID: 0000-0003-2310-9415}
\quad$\vert$\quad
\href{https://github.com/VigneshAstro}{GitHub}}
\vspace{5pt}
\color{rulegray}\hrule\color{black}
\vspace{2pt}
% ===============================================================
\section{Academic Appointments}

\grantheading
{Specially Appointed Assistant Professor}
{June 2026--present}
\emph{Astronomical Institute, Graduate School of Science,
Tohoku University; Sendai, Japan}

{\color{entrylight}
\vspace{-1.0mm}Research on exoplanet atmospheres, planetary-system architecture,
and high-resolution spectroscopy.
\par}

%\vspace{4pt}

\grantheading
{Visiting Researcher}
{January--April 2026}
\emph{Astrobiology Center, NINS;
Tokyo, Japan}

{\color{entrylight}
\vspace{-1.0mm}Instrumentation development and observations with a new
high-resolution spectrograph.
\par}

%\vspace{4pt}

\grantheading
{Postdoctoral Researcher}
{Nov 2022--April 2026}
\emph{Trottier Space Institute, McGill University;
Montreal, Canada}

{\color{entrylight}
\vspace{-1.0mm}Atmospheric characterization of transiting exoplanets using
ground- and space-based telescopes.
\par}

%\vspace{4pt}


\grantheading
{Project Research Staff}
{Oct 2021--Sept 2022}
\emph{Exoplanet Search Project Office, Astrobiology Center,
NINS; Tokyo, Japan}

{\color{entrylight}
\vspace{-1.0mm}Atmospheric characterization of young transiting planets using the
InfraRed Doppler spectrograph on the Subaru Telescope.
\par}



% ===============================================================
\section{Professional Development}

\grantheading
{Mini-MBA: Executive Development Course}
{October 2025}
\emph{Desautels Faculty of Management, McGill University}

{\color{entrylight}
\vspace{-1.0mm}Intensive executive training in leadership, strategy, and innovation
for research and technology professionals.
\par}

%\vspace{4pt}


\grantheading
{International Diploma in Astrobiology}
{March 2016}
\emph{Indian Astrobiology Research Centre; Mumbai, India}

{\color{entrylight}
\vspace{-1.0mm}IARC International Diploma Program; Excellence Grade: A.
\par}



% ===============================================================
\section{Education}

\grantheading
{Doctoral Research in Exoplanet Science}
{September 2021}
\emph{Tokyo Institute of Technology; Tokyo, Japan}

{\color{entrylight}
\vspace{-1.0mm}Thesis: \emph{Atmospheric Evolution of Low-Mass-Star Planets:
A Tale of Evolution}.

\vspace{-1.0mm}Recipient of the Monbukagakusho (MEXT) Scholarship from the
Government of Japan.
\par}

%\vspace{4pt}


\grantheading
{Master's Degree in Aerospace Engineering}
{May 2017}
\emph{Indian Institute of Technology Kanpur, India;
GPA: 9.0/10.0}

{\color{entrylight}
\vspace{-1.0mm}Recipient of the President Dr.\ Shankar Dayal Sharma Gold Medal.
\par}

%\vspace{4pt}


\grantheading
{Bachelor's Degree in Aerospace Engineering}
{June 2015}
\emph{Amrita Vishwa Vidyapeetham; Coimbatore, India;
GPA: 8.58/10.0}
% ===============================================================
\section{Honors and Achievements}

\begingroup
\setlength{\tabcolsep}{0pt}
\renewcommand{\arraystretch}{0.216}

\begin{tabularx}{\textwidth}{
    @{}
    >{\vspace{0pt}\raggedleft\arraybackslash
      \bfseries\color{slate}}p{1.2cm}
    @{\hspace{0.5cm}}
    >{\vspace{0pt}\raggedright\arraybackslash}X
    @{}
}

2026 &
\textbf{Canadian Space Agency JWST Cycle 4 Grant}, Montreal, Canada.
Grant to study a multi-planet system with JWST. \\

2025 &
\textbf{Arena FIDE Master Title (AFM) in Chess}, Montreal, Canada.
Achieved the Arena FIDE Master title in early 2025. \\

2024 &
\textbf{Private Pilot Training}, Montreal, Canada. \\

2023 &
\textbf{Canadian Transnational U1300 Chess Champion}, Montreal, Canada. \\

2020 &
Graduate Student Assistant--Program Developer (\textbf{GSA-D}),
Tokyo Tech, Japan. \\

2018 &
\textbf{Recipient of the Monbukagakusho (MEXT) Scholarship},
Government of Japan, Tokyo Institute of Technology, Japan. \\

2017 &
\textbf{President Shankar Dayal Sharma Gold Medal},
Indian Institute of Technology Kanpur, India.
Highest honor for graduate studies at IIT Kanpur. \\

2016 &
Published the novel \textit{The Forest Whisperings}
with Partridge Publishing, Coimbatore, India. \\

2016 &
\textbf{Best Paper Award for Motivational Research},
ICMME 2016, Shanghai, China. \\

2015 &
\textbf{Government of India Assistantship for Graduate Studies},
Indian Institute of Technology Kanpur, India. \\

2009 &
\textbf{Honorable Mention},
NASA Space Settlement Contest, NASA Ames Center, USA.

\end{tabularx}
\endgroup

% ===============================================================
\section{Publications}

\textbf{Primary-author publications (first or second author):}

\begin{itemize}

\item
\textbf{Krishnamurthy, V.}, Cowan, N. B., Fukagawa, M., \& Hirano, T. (2026),
``Beyond five scale heights: Composition- and cloud-dependent Ariel Tier-2
requirements for sub-Neptunes.''
\href{https://arxiv.org/abs/2609.01577}{arXiv:2609.01577}
[\emph{under review}].

\item
\textbf{Krishnamurthy, V.} (2026),
``Energetic decoupling from biological scaling: A dimensionless framework
for technological civilizations.''
[\emph{under review in Acta Astronautica}].

\item
\textbf{Krishnamurthy, V.}, Carteret, Y., et al. (2026),
``Continuous helium absorption from both the leading and trailing tails of WASP-107 b.''
\href{https://doi.org/10.1038/s41550-025-02710-8}
{\emph{Nature Astronomy}, 10, 258--270}.

\item
\textbf{Krishnamurthy, V.} \& Cowan, N. B. (2024),
``Helium in Exoplanet Exospheres: Orbital and Stellar Influences.''
\href{https://doi.org/10.3847/1538-3881/ad5441}
{\emph{The Astronomical Journal}, 168, 30}.

\item
\textbf{Krishnamurthy, V.}, Hirano, T., et al. (2023),
``Absence of extended atmospheres in low-mass star radius-gap planets GJ 9827 b, GJ 9827 d, and TOI-1235 b.''
\href{https://doi.org/10.1093/mnras/stad404}
{\emph{Monthly Notices of the Royal Astronomical Society}, 521, 1210--1220}.

\item
\textbf{Krishnamurthy, V.}, Hirano, T., Stefansson, G., et al. (2021),
``Non-detection of helium in the upper atmospheres of TRAPPIST-1b, e, and f.''
\href{https://doi.org/10.3847/1538-3881/ac0d57}
{\emph{The Astronomical Journal}, 162, 82}.

\item
Hirano, T., \textbf{Krishnamurthy, V.}, Gaidos, E., et al. (2020),
``Limits on the spin--orbit angle and atmospheric escape for the
22-Myr-old planet AU Mic b.''
\href{https://doi.org/10.3847/2041-8213/aba6eb}
{\emph{The Astrophysical Journal Letters}, 899, L13}.

\item
\textbf{Krishnamurthy, V.}, et al. (2015),
``Computational study of the aerodynamics of the gliding snake
\textit{Chrysopelea paradisi}.''
\href{https://www.academia.edu/69469489/Computational_study_of_the_aerodynamics_of_the_gliding_snake_Chrysopelea_Paradisi}
{\emph{International Journal of Applied Engineering Research},
10, 14476--14479}.

\end{itemize}


\textbf{Invited summary and research briefing:}

\begin{itemize}

\item
\textbf{Krishnamurthy, V.} (2026),
``JWST tracks helium escape from the atmosphere of the exoplanet WASP-107b.''
\href{https://doi.org/10.1038/s41550-025-02722-4}
{\emph{Nature Astronomy}, 10, 185--186 (Research Briefing)}.

\end{itemize}


\textbf{Co-authored publications:}

\begin{itemize}

\item
\begin{sloppypar}
L'Heureux, A., et al., including \textbf{Krishnamurthy, V.} (2026),
``Radial velocity detection of the TRAPPIST-\allowbreak 1 planetary system with
SPIRou and NIRPS.''
\href{https://doi.org/10.1051/0004-6361/202659920}
{\emph{Astronomy \& Astrophysics}, 714, A359}.
\end{sloppypar}

\item
Srivastava, A., et al., including \textbf{Krishnamurthy, V.} (2026),
``Night-sky emission correction techniques for high-resolution spectroscopy:
Demonstration with NIRPS.''
\href{https://doi.org/10.1051/0004-6361/202556093}
{\emph{Astronomy \& Astrophysics}, 714, A123}.

\item
Kawauchi, K., et al., including \textbf{Krishnamurthy, V.} (2026),
``A search for helium in the atmospheres of three sub-Neptunes and a
super-Earth around M-dwarfs.''
\href{https://doi.org/10.1093/pasj/psag104}
{\emph{Publications of the Astronomical Society of Japan},
78, 1815--1828}.

\item
Wardenier, J. P., et al., including \textbf{Krishnamurthy, V.} (2026),
``Parametrizing the projected wind fields of ultra-hot Jupiters in
thermal emission: An application to GCM spectra of WASP-76b.''
\href{https://doi.org/10.1051/0004-6361/202659973}
{\emph{Astronomy \& Astrophysics}, 713, A146}.

\item
Pelletier, S., et al., including \textbf{Krishnamurthy, V.} (2026),
``The breaking of cold traps and onset of titanium in ultra-hot Jupiter
atmospheres: WASP-189b in context.''
\href{https://doi.org/10.1051/0004-6361/202659885}
{\emph{Astronomy \& Astrophysics}, 711, A119}.

\item
Weisserman, D., et al., including \textbf{Krishnamurthy, V.} (2026),
``Super-Earth masses and stellar abundances from NIRPS reveal tentative
evidence for water-rich formation around M dwarfs.''
\href{https://doi.org/10.1051/0004-6361/202558011}
{\emph{Astronomy \& Astrophysics}, 709, A165}.

\item
Fukuda, I., et al., including \textbf{Krishnamurthy, V.} (2026),
``The mass of TOI-1883 b: A low-density super-Neptune in the ridge regime
transiting an early-M dwarf.''
\href{https://doi.org/10.1093/pasj/psag084}
{\emph{Publications of the Astronomical Society of Japan}, 78, 1602--1614}.

\item
Lamontagne, P., et al., including \textbf{Krishnamurthy, V.} (2026),
``NIRPS tightens the mass estimate of GJ 3090 b and detects a planet near
the stellar rotation period.''
\href{https://doi.org/10.1051/0004-6361/202556887}
{\emph{Astronomy \& Astrophysics}, 706, A278}.

\item
Lacedelli, G., et al., including \textbf{Krishnamurthy, V.} (2026),
``Sibling sub-Neptunes around sibling M dwarfs: TOI-521 and TOI-912.''
\href{https://doi.org/10.1051/0004-6361/202557304}
{\emph{Astronomy \& Astrophysics}, 705, A260}.

\item
Allart, R., et al., including \textbf{Krishnamurthy, V.} (2025),
``Complex structure of escaping helium spanning more than half the orbit
of the ultra-hot Jupiter WASP-121b.''
\href{https://doi.org/10.1038/s41467-025-66628-5}
{\emph{Nature Communications}, 16, 10822}.

\item
Parc, L., et al., including \textbf{Krishnamurthy, V.} (2025),
``NIRPS and TESS reveal a peculiar system around the M dwarf TOI-756:
A transiting sub-Neptune and a cold eccentric giant.''
\href{https://doi.org/10.1051/0004-6361/202555684}
{\emph{Astronomy \& Astrophysics}, 702, A138}.

\item
Ikuta, K., et al., including \textbf{Krishnamurthy, V.} (2025),
``The mass of TOI-654b: A short-period sub-Neptune transiting a
mid-M dwarf.''
\href{https://doi.org/10.1093/pasj/psaf090}
{\emph{Publications of the Astronomical Society of Japan}, 77, 1101--1112}.

\item
Splinter, J., et al., including \textbf{Krishnamurthy, V.} (2025),
``Precise constraints on the energy budget of WASP-121b from its
JWST NIRISS/SOSS phase curve.''
\href{https://doi.org/10.3847/1538-3881/ae0e52}
{\emph{The Astronomical Journal}, 170, 323}.

\item
Bazinet, L., et al., including \textbf{Krishnamurthy, V.} (2025),
``Quantifying thermal water dissociation in the dayside photosphere
of WASP-121b using NIRPS.''
\href{https://doi.org/10.1051/0004-6361/202553724}
{\emph{Astronomy \& Astrophysics}, 701, A276}.

\item
Gomes da Silva, J., et al., including \textbf{Krishnamurthy, V.} (2025),
``Blind search for activity-sensitive lines in the near-infrared using
HARPS and NIRPS observations of Proxima and Gl 581.''
\href{https://doi.org/10.1051/0004-6361/202555013}
{\emph{Astronomy \& Astrophysics}, 700, A177}.

\item
Suárez Mascareño, A., et al., including
\textbf{Krishnamurthy, V.} (2025),
``Diving into the planetary system of Proxima with NIRPS:
Breaking the metre-per-second barrier in the infrared.''
\href{https://doi.org/10.1051/0004-6361/202553728}
{\emph{Astronomy \& Astrophysics}, 700, A11}.

\item
Bouchy, F., et al., including \textbf{Krishnamurthy, V.} (2025),
``NIRPS joining HARPS at the ESO 3.6-m telescope:
On-sky performance and science objectives.''
\href{https://doi.org/10.1051/0004-6361/202453341}
{\emph{Astronomy \& Astrophysics}, 700, A10}.

\item
Vaulato, V., et al., including \textbf{Krishnamurthy, V.} (2025),
``Hydride-ion continuum hides absorption signatures in the NIRPS
near-infrared transmission spectrum of the ultra-hot gas giant WASP-189b.''
\href{https://doi.org/10.1051/0004-6361/202452972}
{\emph{Astronomy \& Astrophysics}, 700, A9}.

\item
Allart, R., et al., including \textbf{Krishnamurthy, V.} (2025),
``NIRPS detection of delayed atmospheric escape from the warm and
misaligned Saturn-mass exoplanet WASP-69b.''
\href{https://doi.org/10.1051/0004-6361/202452525}
{\emph{Astronomy \& Astrophysics}, 700, A7}.

\item
Artigau, É., et al., including \textbf{Krishnamurthy, V.} (2024),
``NIRPS first light and early science: Breaking the 1 m\,s$^{-1}$
radial-velocity precision barrier at infrared wavelengths.''
\href{https://doi.org/10.1117/12.3018994}
{\emph{Ground-based and Airborne Instrumentation for Astronomy X},
Proc. SPIE 13096, 130960C}.

\item
Kuzuhara, M., et al., including \textbf{Krishnamurthy, V.} (2024),
``Gliese 12b: A temperate Earth-sized planet at 12 pc ideal for
atmospheric transmission spectroscopy.''
\href{https://doi.org/10.3847/2041-8213/ad3642}
{\emph{The Astrophysical Journal Letters}, 967, L21}.

\item
Kagetani, T., et al., including \textbf{Krishnamurthy, V.} (2023),
``The mass of TOI-519 b: A close-in giant planet transiting a
metal-rich mid-M dwarf.''
\href{https://doi.org/10.1093/pasj/psad031}
{\emph{Publications of the Astronomical Society of Japan}, 75, 713--721}.

\item
Hirano, T., et al., including \textbf{Krishnamurthy, V.} (2023),
``An Earth-sized planet around an M5 dwarf star at 22 pc.''
\href{https://doi.org/10.3847/1538-3881/acb7e1}
{\emph{The Astronomical Journal}, 165, 131}.

\item
Morello, G., et al., including \textbf{Krishnamurthy, V.} (2023),
``TOI-1442b and TOI-2445b: Two potentially rocky ultra-short period planets
around M dwarfs.''
\href{https://doi.org/10.1051/0004-6361/202243592}
{\emph{Astronomy \& Astrophysics}, 673, A32}.

\item
Delrez, L., et al., including \textbf{Krishnamurthy, V.} (2022),
``Two temperate super-Earths transiting a nearby late M dwarf.''
\href{https://doi.org/10.1051/0004-6361/202244041}
{\emph{Astronomy \& Astrophysics}, 667, A59}.

\item
Mori, M., et al., including \textbf{Krishnamurthy, V.} (2022),
``TOI-1696: A nearby M4 dwarf with a 3\,R$_{\oplus}$ planet in the
Neptunian desert.''
\href{https://doi.org/10.3847/1538-3881/ac6bf8}
{\emph{The Astronomical Journal}, 163, 298}.

\item
Kawauchi, K., et al., including \textbf{Krishnamurthy, V.} (2022),
``Validation and atmospheric exploration of the sub-Neptune TOI-2136b
around a nearby M3 dwarf.''
\href{https://doi.org/10.1051/0004-6361/202243381}
{\emph{Astronomy \& Astrophysics}, 666, A4}.

\item
Harakawa, H., et al., including \textbf{Krishnamurthy, V.} (2022),
``A super-Earth orbiting near the inner edge of the habitable zone
around the M4.5 dwarf Ross 508.''
\href{https://doi.org/10.1093/pasj/psac044}
{\emph{Publications of the Astronomical Society of Japan}, 74, 904--922}.

\item
Gaidos, E., et al., including \textbf{Krishnamurthy, V.} (2020),
``Zodiacal exoplanets in time XI: The orbit and radiation environment
of the young M-dwarf-hosted planet K2-25b.''
\href{https://doi.org/10.1093/mnrasl/slaa136}
{\emph{Monthly Notices of the Royal Astronomical Society: Letters},
498, L119--L124}.

\end{itemize}

% ===============================================================
\section{Conference Presentations}

\presentationentry
{\href{https://nibb-abc-conf70.nibb.ac.jp/home}{70th NIBB Conference / ABC International Symposium: Bridging Astronomy and Biology: Spectral Footprints of Life Beyond Earth} (oral presentation)}
{Okazaki, Japan}
{Energetic Decoupling from Biological Scaling: A Dimensionless Framework for Technological Civilizations}
{11 November 2026}

\presentationentry
{HWO Japan Science Workshop 2 (presentation)}
{JAXA, Sagamihara}
{Atmosphere and escape in the HWO era: What JWST is teaching us}
{July 2026}

\presentationentry
{Exoplanets VI (presentation)}
{Porto, Portugal}
{Know thy Star, know thy Planet: Probing TLSE in a M-star}
{June-July 2026}

\presentationentry
{Cool Stars 23 (poster)}
{Tokyo, Japan}
{Exploring TLSE in a M-star}
{June 2026}

\presentationentry
{NIRPS Science Meeting (presentation)}
{Geneva, Switzerland}
{Young planets with escaping atmospheres}
{November 2025}

\presentationentry
{CRAQ Annual Meeting 2025 (presentation)}
{Quebec, Canada}
{NIRISS/SOSS observations of WASP-107b}
{May 2025}

\presentationentry
{Two HoRSEs (poster)}
{Berlin, Germany}
{Atmosphere of the young Neptune DS Tuc Ab}
{July 2024}

\presentationentry
{CRAQ Annual Meeting 2024 (flash presentation)}
{Quebec, Canada}
{Atmospheres of young planets}
{May 2024}

\presentationentry
{NIRPS Science Meeting (presentation)}
{Montreal, Canada}
{Atmospheres of young planets}
{April 2024}

\presentationentry
{NIRPS Science Meeting (presentation)}
{Montreal, Canada}
{Search for helium in NIRPS targets}
{October 2023}

\presentationentry
{Exoclimes VI (poster)}
{Exeter, UK}
{Near-infrared observations of transiting-planet atmospheres}
{June 2023}

\presentationentry
{CRAQ Annual Meeting 2023 (presentation)}
{Quebec, Canada}
{Atmospheres of radius-gap planets}
{May 2023}

\presentationentry
{Japan Geoscience Union Meeting 2022 (presentation)}
{Online}
{Atmospheres of low-mass-star planets: Atmospheric escape and beyond}
{May 2022}

\presentationentry
{Exoplanets IV, AASTCS 9 (poster)}
{Las Vegas, USA}
{Near-infrared observations of TOI-1235b, GJ 9827b, and GJ 9827d}
{May 2022}

\presentationentry
{53\textsuperscript{rd} Annual Meeting of the Division for Planetary Sciences (iPoster)}
{Online}
{High-resolution near-infrared observations of the TRAPPIST-1 planets}
{October 2021}

\presentationentry
{Subaru Users Meeting FY2020 (poster)}
{Online}
{Characterization of sub-Neptune and rocky-planet atmospheres with near-infrared spectroscopy}
{March 2021}

\presentationentry
{Rocky Planets in the Era of JWST: Theory and Observation (poster)}
{NASA GSFC, USA}
{Exploring the infrared helium line for atmospheric characterization}
{November 2019}

\presentationentry
{Workshop on TESS Follow-up Obs. for Exoplanet (presentation)}
{Busan, South Korea}
{New observational results from IRD spectrograph on Subaru Telescope}
{October 2019}

\presentationentry
{Tokyo Planetary Seminar (presentation)}
{JAXA, Tanegashima}
{New observational results from IRD spectrograph on Subaru Telescope}
{September 2019}

\presentationentry
{TIEC Research and Presentation 2019--2020 (presentation)}
{Tokyo, Japan}
{Ground-based observations of terrestrial exoplanet systems}
{April 2019}

\presentationentry
{TIEC Research and Presentation 2018--2019 (poster)}
{Tokyo, Japan}
{Current developments in ground-based exoplanet observations}
{November 2018}

\presentationentry
{AIAA Oklahoma Symposium (presentation)}
{Edmond, OK, USA}
{Metamaterial-inspired smart composites for space applications}
{April 2018}

\presentationentry
{6\textsuperscript{th} ELSI International Symposium (poster)}
{Tokyo, Japan}
{Concept of a sustainable habitat on Mars}
{January 2018}

\presentationentry
{Master's Thesis Defense (presentation)}
{Kanpur, India}
{Modeling the elasto-plastic behavior of aluminum alloys using micro- and damage mechanics}
{May 2017}

\presentationentry
{3\textsuperscript{rd} International Conference on Mechanical Engineering (presentation)}
{Shanghai, China}
{MC-DSMC solver for an optimized space-crew re-entry orbital vehicle}
{October 2016}

\presentationentry
{Bachelor's Thesis Defense (presentation)}
{Coimbatore, India}
{Monte Carlo simulations for re-entry vehicles}
{June 2015}


% ===============================================================
\section{Invited Talks and Colloquia}

\invitedtalkentry
{Astronomy Colloquium, Tohoku University}
{Sendai, Japan}
{Dr.\ Misato Fukagawa}
{July 2026}

\invitedtalkentry
{Sato Laboratory, Institute of Science Tokyo}
{Tokyo, Japan}
{Dr.\ Bun'ei Sato}
{April 2026}

\invitedtalkentry
{Young Transiting Planet Workshop 2026}
{Okinawa, Japan}
{Dr.\ Norio Narita}
{March 2026}

\invitedtalkentry
{Astrobiology Center, National Astronomical Observatory of Japan}
{Tokyo, Japan}
{Dr.\ Teruyuki Hirano}
{February 2026}

\invitedtalkentry
{Department of Astronomy \& Astrophysics, TIFR}
{Mumbai, India}
{Dr.\ Joe P. Ninan}
{December 2025}

\invitedtalkentry
{Inter-University Centre for Astronomy and Astrophysics}
{Pune, India}
{Dr.\ R. Srianand and Dr.\ Anupam Bhardwaj}
{December 2025}

\invitedtalkentry
{Department of SPASE, IIT Kanpur}
{Kanpur, India}
{Dr.\ Prashant Pathak}
{August 2025}

\invitedtalkentry
{Department of Earth \& Space Sciences, IIST}
{Thiruvandrum, India}
{Dr.\ Anand Narayanan}
{July 2025}

\invitedtalkentry
{Center for Astrophysics | Harvard \& Smithsonian}
{Boston, USA}
{Dr.\ Dave Charbonneau}
{October 2024}

\invitedtalkentry
{RIKEN}
{Saitama, Japan}
{Dr.\ Nami Sakai and Dr.\ Yui Kawashima}
{January 2024}

\invitedtalkentry
{Amrita Vishwa Vidyapeetham}
{Coimbatore, India}
{Dr.\ S. Mahadevan and Dr.\ Bharat Kishore Sharma}
{January 2024}

\invitedtalkentry
{Indian Institute of Astrophysics}
{Bengaluru, India}
{Dr.\ Sujan Sengupta and Dr.\ Piyali Chatterjee}
{January 2024}

\invitedtalkentry
{Astrobiology Center, National Astronomical Observatory of Japan}
{Tokyo, Japan}
{Dr.\ Teruyuki Hirano}
{May 2022}

% ===============================================================
\section{Grants and Awards}

\grantheading
{Canadian Space Agency -- JWST Cycle 4 Grant}
{2025--present}
Project: Probing Stellar Behavior and Planetary Diversity in K2-384 \vspace{-1.5mm}
\begin{lightitems}
\item
Canadian Space Agency grant supporting our JWST DDT program in Cycle 4.
\item
The CAD 60,000 grant supports in the publications and conference participation.
\end{lightitems}


\grantheading
{Clark Science Executive Leadership Fund (Clark-SELF)}
{2025}
McGill University \vspace{-1.5mm}
\begin{lightitems}
\item
Recipient of the Clark-SELF award for executive MBA-level leadership
training.
\item
The CAD 7,000 award supported participation in the Mini-MBA Executive
Development Program.
\end{lightitems}


\grantheading
{NASA SEEC Grant}
{2020--2022}
Role: Collaborator; PI: Ravi Kopparapu, NASA Goddard Space Flight Center, USA \vspace{-1.5mm}
\begin{lightitems}
\item
Simulated transmission spectra of atomic helium as a tracer of
atmospheric escape.
\item
The USD 100,000 grant supported the development of helium models.
\end{lightitems}


\grantheading
{Monbukagakusho (MEXT) Scholarship for Doctoral Studies}
{2018--2021}
Ministry of Education, Culture, Sports, Science and Technology,
Government of Japan \vspace{-1.5mm}
\begin{lightitems}
\item
Scholarship supporting doctoral studies in exoplanet atmospheres at
the Tokyo Tech, Japan.
\item
The scholarship provided JPY 6,090,000 over 3.5 years.
\end{lightitems}


\grantheading
{GATE Stipend for Master's Studies}
{2015--2017}
All India Council for Technical Education scholarship,
Government of India \vspace{-1.5mm}
\begin{lightitems}
\item
Scholarship supporting master's studies in aerospace engineering at
the IIT Kanpur, India.
\item
The scholarship provided INR 300,000 over two years.

\end{lightitems}


% ===============================================================

\section{Current and Past Projects}

\projectentry
{Atmospheric and Orbital Architecture of Close-in Packed Planetary Systems}
{2023--present}
{\vspace{-1mm}Role: Principal Investigator}
{\vspace{-1.5mm}Investigating the atmospheres and orbital architectures of compact,
close-in multiplanet systems. Systems such as TRAPPIST-1 and K2-384,
which contain several planets in or near orbital resonances, provide
valuable insights into planetary formation and evolutionary pathways.}


\projectentry
{NIRPS -- Search for Blue Planets Around Red Stars}
{2022--present}
{\vspace{-1mm}Role: Instrument and Guaranteed Time Observations team member}
{\vspace{-1.5mm}Member of the NIRPS instrument team and observer since 2022.
The project focuses on transit observations of young planets orbiting
red dwarfs using the Near Infra-Red Planet Searcher on the ESO 3.6-m
telescope at La Silla Observatory, Chile.}


\projectentry
{Subaru Strategic Program -- InfraRed Doppler Spectrograph}
{2019--2024}
{\vspace{-1mm}Role: Co-investigator; PI: Bun'ei Sato}
{\vspace{-1.5mm}Member of the Subaru Strategic Program team since 2019. The program
focuses on discovering Earth-mass planets in the habitable zones of
mid-to-late M dwarfs and determining the distribution of planets around
cool stars. The program comprises 175 observing nights over five years
with the high-resolution IRD spectrograph on the 8.2-m Subaru Telescope.
My work focuses on stellar characterization using the He\,I triplet at
1083~nm.}


\projectentry
{Simulating Transit Spectra of Atomic Helium as a Tracer of Escape}
{2020--2022}
{\vspace{-1mm}Role: Collaborator; PI: Ravi Kopparapu}
{\vspace{-1.5mm}The project updated a 3D upper-atmosphere model to
calculate atmospheric escape rates and predict transmission spectra for
transiting exoplanets around a range of stellar types. The model
incorporated atomic absorption lines, particularly the metastable
helium triplet at 10833~\AA, which traces the upper-atmospheric region
from approximately $\mu$bar to nbar where escape
is launched.}


\projectentry
{IRD General Observing Programs}
{2019--2022}
{\vspace{-1mm}Role: Principal Investigator and Co-investigator}
{\vspace{-1.5mm}Secured more than 20 hours of Subaru Telescope time as Principal
Investigator, in addition to several observing nights as a
Co-investigator, using the IRD spectrograph to investigate
photoevaporation in sub-Neptunes. The programs resulted in more than
three peer-reviewed publications.}


% ===============================================================
\section{Research Supervision}

\grantheading
{Graduate Supervision, IIST}
{2025--2026}
\emph{Co-supervisor for Arya Raj with Prof.\ Anand Narayanan and Prof.\ Mathew}

Project: Feasibility of detecting oxygen with the Thirty Meter Telescope
\begin{lightitems}
\item
\vspace{-1.5mm}Simulated the atmosphere of a habitable-zone Earth-like planet orbiting
a Sun-like star.
\item
Assessed oxygen detectability using high-resolution cross-correlation
spectroscopy.
\end{lightitems}


\grantheading
{Graduate Supervision, Université de Montréal}
{2024--2026}
\emph{Co-supervisor for Mathis Bouffard with Prof.\ René Doyon and
Prof.\ Nicolas Cowan}

Project: High-resolution transit spectroscopy of the TRAPPIST-1 planets
\begin{lightitems}
\item
\vspace{-1.5mm}Expanded Mathis's undergraduate research to include the
complete TRAPPIST-1 planetary system.
\item
Analyzed transit observations obtained with CFHT/SPIRou and Subaru/IRD.
\end{lightitems}


\grantheading
{Undergraduate Supervision, University of British Columbia}
{2025--2026}
\emph{Co-supervisor for Sean Collins with Prof.\ Michelle Kunimoto}

Project: Confirmation of a TESS planet candidate
\begin{lightitems}
\item
\vspace{-1.5mm}Analyzed TESS observations and ground-based photometry to confirm a
TESS Object of Interest.
\item
Supported the preparation of radial-velocity observations with a
high-resolution spectrograph.
\end{lightitems}


\grantheading
{Undergraduate Supervision, Amrita Vishwa Vidyapeetham}
{2024--2026}
\emph{Supervisor for Anusha and Ashfaq Ali (non-thesis research)}

Project: Transit-timing variation analysis of a multiplanet M-dwarf system
\begin{lightitems}
\item
\vspace{-1.5mm}Investigated a multiplanet system around a M dwarf to constrain
planetary masses using TTVs.
\item
Combined Kepler, TESS, and ground-based photometry to measure the
transit-timing variations.
\end{lightitems}


\grantheading
{Undergraduate Honors Research Thesis, McGill University}
{2023--2024}
\emph{Co-supervisor for Mathis Bouffard and Maddy Walkington with
Prof.\ Nicolas Cowan}

Project: High-resolution transit spectroscopy of TRAPPIST-1b
\begin{lightitems}
\item
\vspace{-1.5mm}Co-supervised the analysis of CFHT/SPIRou transit spectroscopy of
TRAPPIST-1b.
\item
The research was presented at the Emerging Researchers in Exoplanet
Science Symposium IX at Cornell University, Ithaca, New York, in
July 2024.
\end{lightitems}

% ===============================================================
\section{Technical and Professional Skills}

\textbf{\color{navy}Software and GitHub Packages}
\hfill
\href{https://github.com/VigneshAstro}{GitHub: VigneshAstro}

\grantheading
{HELIX -- Helium Exoplanet Literature Investigation Xtractor}
{2024}
\emph{Python wrapper for extracting helium-related papers from arXiv}
\begin{lightitems}
\item
\vspace{-1.5mm}Developed HELIX, a Python wrapper for efficiently extracting research papers from arXiv.
\item
Created an accompanying Jupyter notebook with instructions for
customizing search terms and extracting literature related to other
atmospheric species, including hydrogen and water.
\end{lightitems}


\grantheading
{Signal-to-Noise Ratio Simulator for NIRPS and HARPS}
{2023}
\emph{Exposure-time and signal-to-noise simulator for HARPS and NIRPS}
\begin{lightitems}
\item
\vspace{-1.5mm}Developed a signal-to-noise ratio simulator for the HARPS and NIRPS
spectrographs on the ESO 3.6-m telescope at La Silla Observatory, Chile.
\item
Improved the ETC and implemented planetary velocity-smearing models for transit observations.
\end{lightitems}


\textbf{\color{navy}Observational Experience}

\grantheading
{Support Astronomer: ESO 3.6-m Telescope}
{2022--present}
\emph{Instruments: NIRPS and HARPS}
\begin{lightitems}
\item
\vspace{-1.5mm}Member of the NIRPS instrument and Guaranteed Time Observations team,
focusing on the atmospheric characterization of transiting exoplanets.
\item
Provided on-site observing support at ESO's La Silla Observatory and
participated in remote observations for the NIRPS GTO program.
\end{lightitems}


\grantheading
{Support Astronomer: Subaru Telescope}
{2020--2022}
\emph{Instrument: InfraRed Doppler spectrograph}
\begin{lightitems}
\item
\vspace{-1.5mm}Worked as a support astronomer for the IRD spectrograph on the Subaru Telescope at Hawaii.
\item
Served as a remote observer for the Subaru Strategic Program from the NAOJ Office in Mitaka.
\end{lightitems}


\textbf{\color{navy}Peer Review}

\grantheading
{Reviewer for Peer-Reviewed Journals}
{2021--present}
\begin{lightitems}
\item
\vspace{-1.5mm}American Astronomical Society journals:
\emph{The AJ},
\emph{The ApJ}, and
\emph{The ApJL}.
\item
\emph{Monthly Notices of the Royal Astronomical Society}.
\item
Nature Portfolio journals, including
\emph{Nature} and \emph{Nature Astronomy}.
\end{lightitems}

% ===============================================================
\section{Professional Service and Committees}

\grantheading
{\href{https://exoclimes.org/}{Exoclimes VII and ExoSlam}}
{2025}
\emph{Local Organizing Committee and Summer School Organizer;
Montreal, Canada}
\begin{lightitems}
\item
\vspace{-1.5mm}Served on the Local Organizing Committee for the seventh edition of
the Exoclimes conference.
\item
Contributed to the organization and delivery of ExoSlam, a summer
school designed for graduate students and early-career researchers.
\end{lightitems}


\grantheading
{Gemini Strategic Science Plan Working Group}
%{September--November 2024}
{2024}
\emph{Exoplanet Representative; online}
\begin{lightitems}
\item
\vspace{-1.5mm}Appointed to the working group responsible for drafting recommendations for the Gemini Board of Directors and shaping the observatory's strategic plan for the 2030s.
\item
Represented the priorities and perspectives of the exoplanet and Canadian astronomical communities.
\end{lightitems}


\grantheading
{\href{https://sites.google.com/view/two-horses/}{Two HoRSEs}}
{July 2024}
\emph{Local Organizing Committee; Berlin, Germany}
\begin{lightitems}
\item
\vspace{-1.5mm}Served on the Local Organizing Committee for the second edition of
\emph{High-Resolution Exoplanet and Stellar Characterization Today and
in the ELT Era}.
\item
Supported the Scientific Organizing Committee with conference
operations, participant coordination, and logistics.
\end{lightitems}


\grantheading
{Trottier Space Institute Seminar Coordinator}
{2023--2024}
\emph{Trottier Space Institute, McGill University; Montreal, Canada}
\begin{lightitems}
\item
\vspace{-1.5mm}Coordinated institute seminars and hosted visiting researchers working
across astronomy, astrophysics, and cosmology.
\end{lightitems}

% ===============================================================
\section{Media and Interviews}

A complete archive of interviews and media coverage is available
\href{https://vigneshastro.github.io/media/} {here}.

\grantheading
{WASP-107b -- Extended Helium Tail}
{2025}
\begin{lightitems}
\item
\vspace{-1.5mm}BBC Tamil:
\href{https://www.bbc.com/tamil/articles/c0jven4j0ngo}
{Extensive helium tail around WASP-107b} and
\href{https://www.bbc.com/tamil/articles/c4gr40qw53xo}
{interview with Dr. Krishnamurthy}.
\item
McGill University:
\href{https://www.mcgill.ca/newsroom/channels/news/exoplanet-observed-shedding-its-atmosphere-real-time-369356}
{Exoplanet observed shedding its atmosphere in real time}.
\item
Trottier Institute for Research on Exoplanets:
\href{https://exoplanetes.umontreal.ca/en/exoplanet-caught-shedding-its-atmosphere-in-real-time/}
{Exoplanet shedding its atmosphere in real time}.
\item
Science Springs:
\href{https://sciencesprings.wordpress.com/2025/12/02/from-the-university-of-chicago-and-mcgill-university-universite-mcgill-ca-and-webb-exoplanet-is-observed-shedding-its-atmosphere-in-real-time/}
{Exoplanet observed shedding its atmosphere in real time}.
\item
Space.com:
\href{https://www.space.com/astronomy/james-webb-space-telescope/giant-helium-cloud-caught-escaping-from-puffy-exoplanet-by-jwst}
{A ``super-puff'' exoplanet is losing its atmosphere, and JWST had a look}.
\item
MyScience:
\href{https://www.myscience.org/en/news/2025/exoplanet_is_observed_shedding_its_atmosphere_in_real_time-2025-mcgill}
{Exoplanet observed shedding its atmosphere in real time}.
\item
Time.news:
\href{https://time.news/helium-leaks-from-exoplanet-wasp-107b-webb-telescope-discovery/}
{Helium leaks from exoplanet WASP-107b}.
\item
Mirage News:
\href{https://www.miragenews.com/exoplanet-spotted-shedding-atmosphere-live-1580797/}
{Exoplanet spotted shedding its atmosphere in real time}.
\item
ScienMag:
\href{https://scienmag.com/helium-absorption-detected-in-wasp-107bs-tails/}
{Helium absorption detected in WASP-107b's tails}.
\end{lightitems}


\grantheading
{Gliese 12b -- Discovery}
{July--August 2024}
\begin{lightitems}
\item
\vspace{-1.5mm}The Conversation:
\href{https://theconversation.com/the-discovery-of-a-new-earth-like-planet-could-shed-further-light-on-what-makes-a-planet-habitable-234127}
{The discovery of a new Earth-like planet could shed further light on
what makes a planet habitable}.
\item
Puthiya Thalaimurai:
\href{https://www.puthiyathalaimurai.com/technology/a-planet-like-earth-has-surprised-scientists}
{Discovery of a potentially habitable-zone exoplanet}.
\item
Puthiya Thalaimurai:
\href{https://youtu.be/AXB-5W64Iac}
{Video report on the discovery of Gliese 12b}.
\item
Trottier Institute for Research on Exoplanets:
\href{https://exoplanetes.umontreal.ca/en/discovery-of-an-earth-like-exoplanet-in-the-habitable-zone-interview-with-vigneshwaran-krishnamurthy/}
{Discovery of an Earth-sized exoplanet in the habitable zone:
Interview with Dr. Krishnamurthy}.
\item
McGill University:
\href{https://www.mcgill.ca/newsroom/channels/news/expert-discovery-new-earth-planet-could-shed-light-what-makes-planet-habitable-358213}
{Discovery of a new Earth-like planet could shed light on what makes
a planet habitable}.
\end{lightitems}

\grantheading
{\href{https://www.puthiyathalaimurai.com/technology/james-webb-telescope-invented-a-new-emerging-solar-family}
{New Star Formation Observed with JWST}}
{November 2023}
\emph{Contributor; Puthiya Thalaimurai; author: Jayashree; language: Tamil}
\begin{lightitems}
\item
\vspace{-1.5mm}Contributed scientific context for an article about the formation of
the protostar HH 212 as observed with JWST.
\end{lightitems}

\grantheading
{\href{https://www.puthiyathalaimurai.com/technology/surprising-discovery-of-3-planets-like-earth-in-trappist1-family}
{TRAPPIST-1 in the Modern Era of Astronomy}}
{November 2023}
\emph{Contributor; Puthiya Thalaimurai; author: Jayashree; language: Tamil}
\begin{lightitems}
\item
\vspace{-1.5mm}Contributed to a review of the TRAPPIST-1 system and current knowledge
of its planetary architecture and atmospheres.
\end{lightitems}

\grantheading
{\href{https://www.puthiyathalaimurai.com/technology/what-causes-solar-storms-in-sun}
{What Is a Solar Storm?}}
{August 2023}
\emph{Contributor; Puthiya Thalaimurai; author: Jayashree; language: Tamil}

\begin{lightitems}
\item
\vspace{-1.5mm}Explained solar energy production, the physical origin of solar storms,
and the role of ISRO's Aditya-L1 spacecraft in monitoring the Sun.
\end{lightitems}


% ===============================================================
\section{Additional Professional Experience and Outreach}


\grantheading
{Market Research Analyst}
{2024--2025}
\emph{GLiN Corporation, Japan}
\begin{lightitems}
\item
\vspace{-1.5mm}Conducted market research for projects involving India and North
America.
\item
Collected and analyzed data and prepared recommendations based on
market trends and client requirements.
\end{lightitems}

\grantheading
{Public Science Outreach}
{2023--2025}
\emph{Trottier Space Institute, Montreal, Canada}
\begin{lightitems}
\item
\vspace{-1.5mm}Contributed to public astronomy and science outreach activities as a
member of the Trottier Space Institute outreach team.
\end{lightitems}

\grantheading
{SHAD Canada Recruiter}
{September 2023--2025}
\emph{Montreal, Canada}
\begin{lightitems}
\item
\vspace{-1.5mm}Promoted experiential education and youth development programs to
high-school students and encouraged participation in SHAD Canada's
summer programs.
\end{lightitems}

\grantheading
{Part-Time Science Teacher}
{March--September 2021}
\emph{RED Program, Musashi High School, Tokyo, Japan}
\begin{lightitems}
\item
\vspace{-1.5mm}Taught a range of science subjects in English to Japanese high-school
students.
\end{lightitems}

\grantheading
{Data Analyst Intern}
{Dec 2020--June 2021}
\emph{Data Visualization Laboratory, Tokyo, Japan}
\begin{lightitems}
\item
\vspace{-1.5mm}Developed visual representations of astronomical data for public
education and science-outreach projects.
\end{lightitems}

\grantheading
{Part-Time English Teacher}
{Nov2020--Sept 2021}
\emph{Tokyo Global Gateway, Tokyo, Japan}
\begin{lightitems}
\item
\vspace{-1.5mm}Taught more than 500 elementary-, middle-, and high-school students
from Tokyo and the surrounding region.
\end{lightitems}

\grantheading
{Teaching Assistant}
{July--September 2020}
\emph{Tokyo Institute of Technology, Tokyo, Japan}
\emph{Instructors: Prof.\ Kumara, Prof.\ Stacey Doremus,
Prof.\ Eri Ota, and Prof.\ Rie Murakami}
\begin{lightitems}
\item
\vspace{-1.5mm}Supported the collaborative Georgia Tech--Tokyo Tech problem-based
learning courses.
\end{lightitems}

\grantheading
{Content Creator -- edX}
{2019--2021}
\emph{Center for Innovative Teaching and Learning,
Tokyo Institute of Technology}
\emph{Supervisors: Prof.\ Jeffrey S. Cross and Prof.\ Nagahama}
\begin{lightitems}
\item
\vspace{-1.5mm}Served as a Graduate Student Assistant--Program Developer and
contributed to the development of online courses for edX.
\end{lightitems}

\grantheading
{President}
{Aug 2019--Aug 2020}
{\color{entrylight}
Indian Students' Association of Tokyo Tech.
\par}

\grantheading
{Distribution Manager}
{Sept 2017--Sept 2019}
{\color{entrylight}
Habitable Press, a publishing imprint focused on Earth, space, and
the future.
\par}

\grantheading
{Research Intern}
{May 2017--April 2018}
{\color{entrylight}
Blue Marble Space Institute of Science, USA.
\par}

\grantheading
{Teaching Assistant}
{January 2016--June 2017}
\emph{Indian Institute of Technology Kanpur, India}

\begin{lightitems}
\item
\vspace{-1.5mm}Conducted tutorial sessions and supported instruction in Engineering
Graphics and Aerospace Structural Analysis~I.
\end{lightitems}


\grantheading
{President -- Astronomy Club}
{June 2013--May 2014}
\emph{Amrita Vishwa Vidyapeetham, India}
\begin{lightitems}
\item
\vspace{-1.5mm}Organized and conducted night-sky observing sessions near Coimbatore,
India.
\end{lightitems}


\grantheading
{Science Outreach Program}
{January 2012--May 2015}
\begin{lightitems}
\item
\vspace{-1.5mm}Encouraged students from multiple districts across Tamil Nadu to
pursue education and careers in science.
\end{lightitems}

\end{document}
